\documentclass[10pt]{report}
\usepackage{ctex}
\usepackage[utf8]{inputenc}
\usepackage{amsmath}
\usepackage{amssymb}
\usepackage{xcolor}
\usepackage{graphicx}
\usepackage[margin=0.1cm]{geometry} % 设置页边距为 0.5cm，尽可能扩大可用区域
\usepackage{setspace}
\renewcommand{\baselinestretch}{1} % 设置行间距为单倍行距
\usepackage{float}
\usepackage{wrapfig}
\usepackage[utf8]{inputenc}
\usepackage[T1]{fontenc}
\usepackage{ctex}
\usepackage{amsmath}
\usepackage{amssymb}
\usepackage{xcolor}
\usepackage{graphicx}
\usepackage{float}
\usepackage{lmodern}
\usepackage{mathrsfs}
\usepackage{lipsum}
\usepackage{array} 
\usepackage{graphicx}
\usepackage{subcaption} % 支持子图
\usepackage{tikz}
\usepackage{braket}
\usepackage{adjustbox}
\usepackage{footmisc}
\usepackage{array}
\usepackage{multirow}
\usepackage[table]{xcolor}
\usepackage{amsmath}
\usepackage{tcolorbox}
\usepackage{caption}
\usepackage{tabularx}
\captionsetup[figure]{labelformat=empty} % 设置图标题的标签格式为空，即不显示“图 X”
\usepackage{graphicx}
\tcbuselibrary{listings,breakable}
\usepackage{xcolor}
\setlength{\parindent}{0pt}
\captionsetup[figure]{skip=0pt}
\usepackage{array}
\begin{document}
	\section{$U_F-I_F$曲线}
	\subsection{$T=297.75K$}
	{
		\resizebox{\textwidth}{!}{ % 宽度撑满文本宽度，高度自动调整
			\begin{tabular}{|c|*{23}{c|}}
				\hline
				序号 &1 & 2 & 3 & 4 & 5 & 6 & 7 & 8 & 9 & 10 & 11 &12 &13&14&15&16&17&18&19&20&21&22&23  \\
				\hline
				$U_F$(V) &0.350 & 0.360 & 0.370 & 0.380 & 0.390 & 0.400 & 0.410 & 0.420 & 0.430 & 0.440 & 0.450 & 0.460 & 0.470 & 0.480 & 0.490 & 0.500 & 0.510 & 0.520 & 0.530 & 0.540 & 0.550 & 0.560 & 0.570  \\
				\hline
				$I_F$($\mu$A) &0.009 & 0.014 & 0.023 & 0.034 & 0.052 & 0.078 & 0.117 & 0.166 & 0.251 & 0.371 & 0.538 & 0.756 & 1.100 & 1.620 & 2.330 & 3.380 & 4.940 & 7.440 & 10.900 & 15.800 & 23.700 & 34.500 & 50.700
				 \\
				\hline
			\end{tabular}
		}
		\begin{figure}[H]
			\centering
			\begin{subfigure}{0.48\textwidth} % 第一个子图，宽度占文本的48%
				\includegraphics[width=\linewidth]{1.1.1.png}
				\caption{$lnI_F=38.59U_F-18.05$}
				\label{fig:sub1}
			\end{subfigure}
			\hfill % 填充水平间距
			\begin{subfigure}{0.48\textwidth} % 第二个子图
				\includegraphics[width=\linewidth]{1.1.2.png}
				\caption{$I_F=1.45\times 10^{-8}e^{38.59U_F}$}
				\label{fig:sub2}
			\end{subfigure}
			\caption{$T=297.75K$}
			\label{fig1:combined}
		\end{figure}
	}
	对于\[I_F = I_s exp(\frac{qV_F}{kT})\]
	我们可以得到 \[I_s = 1.45\times 10^{-8}\mu A,\frac{q}{kT} = 38.59V^{-1}\]
	我们可以计算出 \[k = \frac{q}{T\times38.59V^{-1}} = \frac{1.6\times10^{-19}C}{297.75K\times38.59V^{-1}}=1.392494325\times10^{-23}J/K\]
	\subsection{$T=312.75K$}
	{
		\resizebox{\textwidth}{!}{ % 宽度撑满文本宽度，高度自动调整
			\begin{tabular}{|c|*{23}{c|}}
				\hline
				序号 &1 & 2 & 3 & 4 & 5 & 6 & 7 & 8 & 9 & 10 & 11 &12 &13&14&15&16&17&18&19&20&21&22&23  \\
				\hline
				$U_F$(V) &0.350 & 0.360 & 0.370 & 0.380 & 0.390 & 0.400 & 0.410 & 0.420 & 0.430 & 0.440 & 0.450 & 0.460 & 0.470 & 0.480 & 0.490 & 0.500 & 0.510 & 0.520 & 0.530 & 0.540 & 0.550 & 0.560 & 0.570  \\
				\hline
				$I_F$($\mu$A) &0.076 & 0.100 & 0.144 & 0.205 & 0.293 & 0.399 & 0.585 & 0.835 & 1.100 & 1.590 & 2.290 & 3.270 & 4.540 & 6.600 & 9.560 & 13.400 & 19.600 & 28.000 & 40.900 & 51.800 & 83.300 & 120.000 & 170.000\\
				\hline
			\end{tabular}
		}
		\begin{figure}[H]
			\centering
			\begin{subfigure}{0.48\textwidth} % 第一个子图，宽度占文本的48%
				\includegraphics[width=\linewidth]{1.2.1.png}
				\caption{$lnI_F=32.9U_F-12.58$}
				\label{fig:sub3}
			\end{subfigure}
			\hfill % 填充水平间距
			\begin{subfigure}{0.48\textwidth} % 第二个子图
				\includegraphics[width=\linewidth]{1.2.2.png}
				\caption{$I_F=3.15\times 10^{-7}e^{35.18U_F}$}
				\label{fig:sub4}
			\end{subfigure}
			\caption{$T=312.75K$}
			\label{fig2:combined}
		\end{figure}
		我们可以得到 \[I_s = 3.15\times 10^{-7}\mu A,\frac{q}{kT} =35.18V^{-1}\]
		我们可以计算出 \[k = \frac{q}{T\times35.18V^{-1}} = \frac{1.6\times10^{-19}C}{312.75K\times35.18V^{-1}}=1.454209003\times10^{-23}J/K\]
	}
	
	\subsection{$T=327.75K$}
	{
		\resizebox{\textwidth}{!}{ % 宽度撑满文本宽度，高度自动调整
			\begin{tabular}{|c|*{23}{c|}}
				\hline
				序号 &1 & 2 & 3 & 4 & 5 & 6 & 7 & 8 & 9 & 10 & 11 &12 &13&14&15&16&17&18&19&20&21&22&23  \\
				\hline
				$U_F$(V) &0.350 & 0.360 & 0.370 & 0.380 & 0.390 & 0.400 & 0.410 & 0.420 & 0.430 & 0.440 & 0.450 & 0.460 & 0.470 & 0.480 & 0.490 & 0.500 & 0.510 & 0.520 & 0.530 & 0.540 & 0.550 & 0.560 & 0.564\\
				\hline
				$I_F$($\mu$A) &0.370 & 0.532 & 0.680 & 0.940 & 1.270 & 1.690 & 2.410 & 3.340 & 4.560 & 6.610 & 8.970 & 12.400 & 17.500 & 24.500 & 34.000 & 47.600 & 66.600 & 93.100 & 129.000 & 183.000 & 261.000 & 362.000 & 578.000\\
				\hline
			\end{tabular}
		}
		\begin{figure}[H]
			\centering
			\begin{subfigure}{0.48\textwidth} % 第一个子图，宽度占文本的48%
				\includegraphics[width=\linewidth]{1.3.1.png}
				\caption{$lnI_F=32.90U_F-12.58$}
				\label{fig:sub5}
			\end{subfigure}
			\hfill % 填充水平间距
			\begin{subfigure}{0.48\textwidth} % 第二个子图
				\includegraphics[width=\linewidth]{1.3.2.png}
				\caption{$I_F=3.44\times 10^{-6}e^{32.90U_F}$}
				\label{fig:sub6}
			\end{subfigure}
			\caption{$T=327.75K$}
			\label{fig3:combined}
		\end{figure}
		我们可以得到 \[I_s = 3.44\times 10^{-6}\mu A,\frac{q}{kT} =32.90V^{-1}\]
		我们可以计算出 \[k = \frac{q}{T\times32.90V^{-1}} = \frac{1.6\times10^{-19}C}{327.75K\times32.90V^{-1}}=1.48382056\times10^{-23}J/K\]
	}
	\begin{figure}[H]
		\centering
		\begin{subfigure}{0.48\textwidth} % 第一个子图，宽度占文本的48%
			\includegraphics[width=\linewidth]{1.1.png}
			\caption{缩略版}
			\label{fig:sub7}
		\end{subfigure}
		\hfill % 填充水平间距
		\begin{subfigure}{0.48\textwidth} % 第二个子图
			\includegraphics[width=\linewidth]{1.2.png}
			\caption{高清版}
			\label{fig:sub8}
		\end{subfigure}
		\caption{三条曲线}
		\label{fig4:combined}
	\end{figure}
	
	
	
	
	
	
	\section{升温曲线}
	{
		\resizebox{\textwidth}{!}{ % 宽度撑满文本宽度，高度自动调整
			\begin{tabular}{|c|*{9}{c|}}
				\hline
				序号 &1 & 2 & 3 & 4 & 5 & 6 & 7 & 8 & 9 \\
				\hline
				$T/K$&301.650 & 303.150 & 313.150 & 323.150 & 333.150 & 343.150 & 353.150 & 363.150 & 373.150\\
				\hline
				$V_F$($V$) &0.564 & 0.558 & 0.543 & 0.510 & 0.480 & 0.455 & 0.430 & 0.406 & 0.382\\
				\hline
			\end{tabular}
		}
		\begin{figure}[H]
			\centering
			\begin{subfigure}{0.48\textwidth} % 第一个子图，宽度占文本的48%
				\includegraphics[width=\linewidth]{2.1.png}
				\caption{摄氏温标:$V_F=-0.00259T+0.64$}
				\label{fig:sub9}
			\end{subfigure}
			\hfill % 填充水平间距
			\begin{subfigure}{0.48\textwidth} % 第二个子图
				\includegraphics[width=\linewidth]{2.2.png}
				\caption{热力学温标:$V_F=-0.00259T+1.35$}
				\label{fig:sub10}
			\end{subfigure}
			\caption{升温曲线}
			\label{fig5combined}
		\end{figure}
		我们可以得到
		\[S = -\frac{k}{e}ln\frac{c}{I_F} = -0.00259V/K\]
		\[U_g = 1.35V\]
		\[E_g = U_g \times 1.6\times10^{-19} C= 2.16\times10^{-19}V\]
	\end{document}